\section{CONCLUSIONS}

This study tested whether a transparent, rule-based expert system could identify at-risk students with performance comparable to standard machine learning baselines. Using data from module BBB, presentation 2013J of the Open University Learning Analytics Dataset, an expert system was built from rules extracted from a trained decision tree and compared against logistic regression and a support vector machine on a held-out test set.

The expert system produced 14 interpretable rules and achieved the highest recall of the three systems, correctly identifying 77 of the 99 at-risk students in the test set. However, the support vector machine achieved the highest accuracy, precision, and F1-score. The results therefore do not show that the expert system was the strongest model overall. Instead, they show that it provided the best recall in this experimental setup while remaining competitive with the ML baselines and offering a clearer decision process.

Average continuous assessment score was the dominant feature in the generated knowledge base, appearing in every rule. Missing or unavailable assessment marks were treated as zero in the weighted assessment calculation, so low values of \texttt{average\_CA\_score} may reflect weak submitted performance, non-submission, or both. Other variables played a smaller secondary role: \texttt{total\_clicks} and \texttt{date\_registration} appeared in specific intermediate score ranges, while \texttt{imd\_band} was the only demographic/background feature retained in the final rules. Other demographic variables, such as gender, region, age band, disability, and highest education, did not contribute to the extracted rule set. This suggests that for this course offering, early assessment performance was a stronger signal of risk than background characteristics.

The main limitation of this study is that it was restricted to a single course offering. Whether the same rules or the same pattern of feature importance would hold across other modules or presentations is not known. Future work could apply the same approach to other OULAD modules to test how well the findings generalise, or involve domain experts in reviewing and refining the generated rules before deployment. The results presented here suggest that rule-based systems can be a practical and transparent alternative when interpretability is a priority, especially when the goal is to make the reasoning behind at-risk flags available to educators.
